\documentclass[a4paper,12pt]{article}
\usepackage[T2A]{fontenc}
\usepackage[utf8]{inputenc}
\usepackage[russian]{babel}
\usepackage{amsmath,amssymb,mathtools}
\usepackage{helvet}
\renewcommand{\familydefault}{\sfdefault}
\usepackage{geometry}
\geometry{left=24mm,right=24mm,top=23mm,bottom=24mm}
\usepackage{microtype,icomma,parskip,setspace}
\usepackage{enumitem,booktabs,tabularx}
\usepackage{tikz}
\usetikzlibrary{arrows.meta,positioning,shapes.geometric,calc}
\usepackage{xcolor,fancyhdr}
\usepackage[hidelinks]{hyperref}
\definecolor{accent}{HTML}{008080}
\definecolor{darkaccent}{HTML}{004D4D}
\definecolor{textgray}{HTML}{333333}
\definecolor{softgray}{HTML}{6B7474}
\color{textgray}
\onehalfspacing
\setlength{\parindent}{0pt}
\setlength{\parskip}{6pt}
\setlist[enumerate]{leftmargin=*,label=\arabic*.,itemsep=9pt,topsep=6pt}
\setlist[itemize]{leftmargin=*,itemsep=3pt,topsep=4pt}
\pagestyle{fancy}
\fancyhf{}
\fancyhead[L]{\small Доли числа и задачи на расход топлива}
\fancyhead[R]{\small Надежда Клименко}
\fancyfoot[C]{\small\thepage}
\renewcommand{\headrulewidth}{0.4pt}
\renewcommand{\footrulewidth}{0pt}
\setlength{\headheight}{15pt}
\makeatletter
\def\headrule{{\color{accent}\hrule\@height\headrulewidth\@width\headwidth \vskip-\headrulewidth}}
\makeatother
\newcommand{\sect}[1]{\section*{#1}\vspace{-3pt}\noindent\textcolor{accent}{\rule{\linewidth}{0.5pt}}\vspace{3pt}}
\newcommand{\key}[1]{\textcolor{darkaccent}{\textbf{#1}}}
\newcommand{\ex}[1]{\medskip\noindent\textcolor{darkaccent}{\textbf{#1}}\par\vspace{-1mm}}
\begin{document}
\thispagestyle{empty}
\begin{center}
\vspace*{28mm}
{\Large\textcolor{darkaccent}{МАТЕМАТИКА}}\par
\vspace{21mm}
{\LARGE Чек-лист}\par
\vspace{14mm}
{\LARGE\textcolor{accent}{Доли числа, целое и части}}\par
\vspace{2mm}
{\LARGE\textcolor{accent}{Задачи на расход топлива}}\par
\vspace{12mm}
{\large Математика · 6 класс · Математическая вертикаль}\par
\vfill
{\large Надежда Клименко}\par
\vspace{12mm}
{\large 31 октября 2026 г.}\par
\vspace{18mm}
{\normalsize Лёвин Артём Александрович}\par
{\normalsize эксклюзивно для Надежды Клименко}\par
\vspace{15mm}
\begin{tikzpicture}
\draw[accent,line width=0.85pt] (0,0) .. controls (1.2,0.65) and (2.0,-0.52) .. (3.1,0)
 .. controls (4.2,0.48) and (5.0,-0.48) .. (6.1,0)
 .. controls (7.2,0.55) and (7.9,-0.38) .. (9.2,0);
\fill[darkaccent] (4.6,0) circle (1.4pt);
\end{tikzpicture}
\end{center}
\clearpage

\sect{1. Три величины в задаче на долю числа}
\key{Целое} $A$ --- число, от которого берут долю. \key{Доля} $f$ --- обыкновенная или десятичная дробь. \key{Часть} $p$ --- величина, соответствующая этой доле.
\[
p=f\cdot A\qquad
A=p:f\ (f\ne0)\qquad
f=p:A\ (A\ne0).
\]
При решении сначала определите, \emph{какое число является целым}, а какое --- его частью.

\ex{Пример 1. Найти часть}
Найдите $\frac13$ от числа $120$:
\[
p=\frac13\cdot120=40.
\]
\ex{Пример 2. Восстановить целое}
Треть числа равна $50$. Чему равно само число?
\[
A=50:\frac13=50\cdot3=150.
\]
\ex{Пример 3. Определить долю}
Из $80$ страниц прочитано $28$. Какая часть книги прочитана?
\[
f=\frac{28}{80}=\frac7{20}.
\]
\key{Сравните формулировки:} «$\frac13$ \emph{от} $120$» требует умножения; «$\frac13$ \emph{неизвестного числа} равна $50$» требует деления на дробь.

\sect{2. Когда целое состоит из нескольких частей}
Если $\frac3{11}$ длины столба находится в земле, то над землёй находится оставшаяся доля:
\[
1-\frac3{11}=\frac8{11}.
\]
\ex{Пример 4. Длина столба двумя способами}
Надземная часть столба равна $5\frac12$ м, а в землю врыто $\frac3{11}$ всей длины. Найдите полную длину $x$.

\key{Первый способ --- по оставшейся доле:}
\[
\frac8{11}x=5\frac12,\qquad
x=\frac{11}{2}:\frac8{11}
=\frac{11}{2}\cdot\frac{11}{8}
=\frac{121}{16}=7\frac9{16}\ \text{м}.
\]
\key{Второй способ --- уравнением по частям:}
\[
\frac3{11}x+5\frac12=x
\quad\Longrightarrow\quad
x-\frac3{11}x=\frac{11}{2}
\quad\Longrightarrow\quad
\frac8{11}x=\frac{11}{2}.
\]
Оба способа дают одинаковый результат. \key{Проверка:} доля надземной части --- $\frac8{11}$, и $\frac8{11}\cdot\frac{121}{16}=\frac{11}{2}$ м.

\sect{3. Расход: литры на один километр}
Пусть $Q$ --- израсходованное топливо (л), $S$ --- пройденный путь (км), $r$ --- расход (л/км). Тогда:
\[
r=\frac QS\qquad
Q=r\cdot S\qquad
S=\frac Qr\ (r>0).
\]
\key{Почему так:} если каждый километр требует $r$ литров, то за $S$ километров потребуется $rS$ литров. Если известен запас $Q$ литров, его хватит на $Q:r$ километров.

\ex{Пример 5. Перевести часть бака в расход}
Объём бака $40$ л. На $400$ км по трассе израсходовали $\frac58$ бака.
\[
Q=40\cdot\frac58=25\ \text{л},\qquad
r_{\text{т}}=\frac{25}{400}=\frac1{16}\ \text{л/км}.
\]
Это означает: на $16$ км расходуется $1$ л бензина, если расход остаётся постоянным.

\sect{4. Что означает «в несколько раз больше»}
Если расход в городе в $k$ раз больше расхода на трассе, то
\[
r_{\text{г}}=k\cdot r_{\text{т}},\qquad k>1.
\]
Например, $1{,}3=\frac{13}{10}$, а $1\frac13=\frac43$. Это \key{разные числа}; при расчётах их нельзя подменять.

\ex{Пример 6. Полный бак для поездки по городу}
В условиях примера 5 расход в городе в $1{,}3$ раза больше, чем на трассе. Тогда
\[
r_{\text{г}}=\frac{13}{10}\cdot\frac1{16}
=\frac{13}{160}\ \text{л/км}.
\]
С полным баком $40$ л можно проехать:
\[
S_{\text{г}}=40:\frac{13}{160}
=40\cdot\frac{160}{13}=\frac{6400}{13}
\approx492{,}3\ \text{км}.
\]
\key{Проверка смысла:} расход в городе больше, значит городская дальность хода меньше трассовой. На трассе полный бак дал бы $40:(1/16)=640$ км.

\key{Обратная проверка:} если потраченные литры вычислить по формуле $Q=rS$, то
\[
Q=\frac{13}{160}\cdot\frac{6400}{13}=40\ \text{л}.
\]
Это исходный объём полного бака, следовательно, вычисления согласованы.

\sect{5. Обыкновенные дроби и здравый смысл}
\begin{itemize}
\item Деление на дробь заменяют умножением на обратную: $a:\frac bc=a\cdot\frac cb$ при $b,c\ne0$.
\item Сокращайте числитель и знаменатель на \emph{одно и то же ненулевое число}: $\frac{12}{15}=\frac45$.
\item Разделяйте расход и дальность хода: $\frac3{40}$ л/км означает $3$ л на $40$ км, то есть $0{,}075$ л/км.
\item Следите за единицами: $\text{л}:\text{км}=\text{л/км}$, а $\text{л}:(\text{л/км})=\text{км}$.
\item Проверяйте результат: для фиксированного запаса более высокий расход должен давать меньший путь.
\end{itemize}

\clearpage
\sect{Алгоритм: как найти целое, долю или часть}
\vspace{8mm}
\begin{center}
\begin{tikzpicture}[
  >=Latex, node distance=14mm,
  start/.style={ellipse,draw=darkaccent,thick,align=center,minimum width=43mm,minimum height=11mm,fill=white,font=\small},
  question/.style={diamond,aspect=2.5,draw=accent,thick,align=center,text width=34mm,inner sep=2pt,font=\small},
  action/.style={rectangle,rounded corners=2pt,draw=darkaccent,thick,align=center,text width=45mm,minimum height=11mm,font=\small},
  arrow/.style={-{Latex[length=3mm]},thick,draw=darkaccent},
  every node/.style={text=textgray}
]
\node[start] (start) {Выделите $A$ --- целое,\\$f$ --- долю, $p$ --- часть};
\node[question,below=17mm of start] (p) {Неизвестна часть $p$?};
\node[action,right=28mm of p] (ap) {$p=f\cdot A$};
\node[question,below=23mm of p] (a) {Неизвестно целое $A$?};
\node[action,right=28mm of a] (aa) {$A=p:f$};
\node[action,below=23mm of a] (af) {Иначе найдите долю:\\$f=p:A$};
\node[start,below=25mm of af] (end) {Подставьте ответ\\и проверьте смысл};
\draw[arrow] (start)--(p);
\draw[arrow] (p.east)--node[above,font=\footnotesize]{да}(ap.west);
\draw[arrow] (p)--node[left,font=\footnotesize]{нет}(a);
\draw[arrow] (a.east)--node[above,font=\footnotesize]{да}(aa.west);
\draw[arrow] (a)--node[left,font=\footnotesize]{нет}(af);
\draw[arrow] (af)--(end);
\draw[arrow] (ap.east)--++(10mm,0)|-(end.east);
\draw[arrow] (aa.east)--++(10mm,0)|-(end.east);
\end{tikzpicture}
\end{center}
\vfill
\noindent\textcolor{softgray}{\small\textbf{Важно.} Для каждого вычисления делением заранее проверьте, что делитель отличен от нуля. В текстовой задаче обязательно укажите единицы результата.}
\clearpage

\sect{Тренировочный блок · Путь А}
\textcolor{softgray}{\small 10 заданий. Выполняйте преобразования последовательно; в текстовых задачах кратко обозначайте целое, часть и долю либо запас, расход и расстояние.}
\vspace{3mm}
\begin{enumerate}
\item Найдите $\dfrac37$ от числа $91$.
\item Какую часть числа $90$ составляет число $24$? Ответ запишите несократимой дробью.
\item Найдите число, если $\dfrac58$ этого числа равны $45$.
\item В землю врыта четверть длины столба, длина которого $16$ м. Сколько метров столба находится над землёй?
\item Над землёй находится $7{,}2$ м столба, а в землю врыто $\dfrac25$ всей его длины. Найдите длину столба.
\item Бак автомобиля вмещает $40$ л. За $250$ км по трассе израсходовали $\dfrac38$ бака. Найдите расход в литрах на километр.
\item По трассе автомобиль расходует $\dfrac2{25}$ л на километр. В городе расход в $1{,}5$ раза больше. Найдите расход по городу.
\item В баке $48$ л топлива. При движении по городу расход составляет $\dfrac3{25}$ л/км. На сколько километров хватит топлива?
\item Бак вмещает $60$ л. На $450$ км трассы израсходовали $\dfrac34$ бака. Расход в городе в $1{,}2$ раза больше трассового. На сколько километров хватит полного бака в городе?
\item Бак вмещает $48$ л. На $360$ км трассы уходит $\dfrac58$ бака. Расход в городе в $1{,}2$ раза больше трассового. На сколько километров хватит полного бака в городе?
\end{enumerate}
\clearpage

\sect{Ответы к тренировочному блоку}
\vspace{6mm}
\begin{center}
\renewcommand{\arraystretch}{1.55}
\begin{tabular}{@{}>{\bfseries}p{18mm}p{113mm}@{}}
\toprule
№ & Ответ \\
\midrule
1 & $39$ \\
2 & $\dfrac4{15}$ \\
3 & $72$ \\
4 & $12$ м \\
5 & $12$ м \\
6 & $\dfrac3{50}$ л/км, или $0{,}06$ л/км \\
7 & $\dfrac3{25}$ л/км, или $0{,}12$ л/км \\
8 & $400$ км \\
9 & $500$ км \\
10 & $480$ км \\
\bottomrule
\end{tabular}
\end{center}
\vspace{9mm}
\sect{Итоговая памятка}
\begin{tabularx}{\linewidth}{@{}p{0.40\linewidth}X@{}}
\toprule
\textbf{Сигнал в условии} & \textbf{Какое действие выполнить} \\
\midrule
«Доля от числа» & Умножьте целое на долю. \\
«Доля числа равна...» & Разделите известную часть на долю. \\
«Какая часть?» & Разделите часть на целое. \\
«В земле $\frac mn$ длины» & Над землёй остаётся $1-\frac mn$ длины. \\
«В $k$ раз больше расход» & Умножьте расход на $k$. \\
«На сколько километров хватит?» & Разделите запас топлива на расход в л/км. \\
\bottomrule
\end{tabularx}
\end{document}
